% Libro 2 · Mundo 7 · Codificación (app: pu-codigos)
\mundo{Codificación}{pu-codigos}{Un código es una regla que transforma un mensaje en otro. El examen da algunos ejemplos
ya codificados y pide aplicar la regla a una palabra nueva, o deducir el significado de un símbolo. Es lógica pura:
comparar, aislar lo que cambia y comprobar.}{%
\item \textbf{7.1} Claves de sustitución
\item \textbf{7.2} Deducir la regla}

\leccion{pu-cd-1}{Claves de sustitución}{Juego de las casillas con fichas que cambian de letra.}

\subsection{La idea: buscar el desplazamiento}
Los códigos más comunes son de \clave{sustitución}: cada letra se cambia por otra según una regla fija. El caso más
famoso es el \clave{cifrado de César}: correr cada letra $k$ lugares en el alfabeto (con $k = 3$: A $\to$ D, B $\to$ E,
…, X $\to$ A). Para descubrir $k$ basta comparar una letra con su código.

\begin{center}\small
\begin{tabular}{c|cccc}
original & M & E & S & A\\
posición & 13 & 5 & 19 & 1\\\hline
código & N & F & T & B\\
posición & 14 & 6 & 20 & 2\\
\end{tabular}
\end{center}

Otras sustituciones: invertir el alfabeto (A $\leftrightarrow$ Z, B $\leftrightarrow$ Y: la posición $p$ pasa a
$27 - p$), cambiar letras por números, o cambiar cada letra por un símbolo.

\subsubsection{Códigos de posición}
A veces la regla no cambia las letras sino su \emph{orden}: invertir la palabra, intercambiar la primera con la última,
mover la primera al final. Se reconoce porque el código usa las mismas letras.

\figura{cesar}{Un código de desplazamiento: cada letra avanza un lugar.}

\begin{regla}
\begin{itemize}
  \item Si el código usa las mismas letras, la regla es de posición. Si usa otras, es de sustitución.
  \item En sustitución, compara las posiciones: ¿suma constante? ¿diferencia constante?
  \item Comprueba la regla con \emph{todas} las letras del ejemplo antes de aplicarla.
\end{itemize}
\end{regla}

\subsection{Ejemplos resueltos}
\begin{ejemplo}[Alfabeto invertido]
Si CASA se escribe XZHZ, ¿cómo se escribe LUNA?
\solucion C (3) $\to$ X (24), A (1) $\to$ Z (26), S (19) $\to$ H (8): suman 27. L (12) $\to$ 15 = O; U (21) $\to$ 6 = F;
N (14) $\to$ 13 = M; A $\to$ Z. LUNA $\to$ OFMZ.
\end{ejemplo}

\begin{ejemplo}[De posición]
Si CAMINO se escribe ONIMAC, ¿cómo se escribe PERRO?
\solucion La palabra se invierte: ORREP.
\end{ejemplo}

\begin{ejemplo}[Números]
Si ROSA se codifica 18-15-19-1, ¿qué palabra es 13-1-18?
\solucion Posiciones: M, A, R. MAR.
\end{ejemplo}

\begin{trampa}
Suponer un desplazamiento con una sola letra cuando la regla podría ser otra: comprueba siempre con la segunda y la
tercera letra.
\end{trampa}

\begin{reto}
En un código, la palabra MESA se escribe NFTB. ¿Cómo se escribe SILLA?
\opciones{TJMMB}{RHKKZ}{TJMMA}{UJNNB}
\end{reto}

\begin{retoresuelto}
Cada letra avanza una posición: M $\to$ N, E $\to$ F, S $\to$ T, A $\to$ B. SILLA $\to$ TJMMB. \textbf{Opción A.}
\end{retoresuelto}

\taller
\piso{Entrena}{codificar}
\begin{ejercicios}
\item Con César $k = 2$, codifica GATO. \sol{ICVQ.}
\item Si 2-5-19-15 significa BESO, ¿qué significa 1-13-15-18? \sol{AMOR.}
\item Si LIBRO se escribe ORBIL, ¿cómo se escribe MESA? \sol{ASEM.}
\end{ejercicios}
\piso{Como en el examen}{opción múltiple}
\begin{ejercicios}
\item Si en un código PATO se escribe QBUP, la palabra que se escribe DBTB es
\opciones{CASA}{EBSA}{CAPA}{COSA}
\sol{A. Cada letra se corre uno hacia atrás.}
\end{ejercicios}

\bitacora

\leccion{pu-cd-2}{Deducir la regla}{Juego de las palabras en otro idioma.}

\subsection{La idea: lo que se repite se corresponde}
Cuando hay varias frases traducidas, se buscan las palabras \clave{comunes}: si dos frases comparten «grande» y
comparten «sapi», entonces «sapi» = «grande». Es como resolver un sistema de ecuaciones: cada frase es una ecuación y
cada coincidencia aísla una incógnita.

\subsubsection{Códigos con símbolos y operaciones}
La misma idea sirve cuando el código es numérico: si \texttt{3 \# 4 = 712} y \texttt{5 \# 2 = 710}, se busca qué
operación da cada parte (aquí, la suma y luego el producto: $3 + 4 = 7$, $3 \cdot 4 = 12$).

\begin{regla}
\begin{itemize}
  \item Compara frases de dos en dos: lo común en ambas lenguas se corresponde.
  \item Descarta lo ya identificado y repite.
  \item Fíjate si el orden de las palabras se conserva (sustantivo y luego adjetivo, o al revés).
\end{itemize}
\end{regla}

\subsection{Ejemplos resueltos}
\begin{ejemplo}[Tres frases]
«ka mo ri» = «río frío grande»; «ri pe» = «grande lago»; «pe ka su» = «lago frío quieto». ¿Qué significa «mo»?
\solucion Las frases 1 y 2 comparten «ri» y comparten «grande»: ri = grande. Entonces, en la 2, pe = lago. Las frases 1 y
3 comparten «ka» y comparten «frío»: ka = frío. En la 1 solo queda mo = río (y en la 3, su = quieto). Fíjate en que el
orden de las palabras no sirvió: la correspondencia salió de las repeticiones.
\end{ejemplo}

\begin{ejemplo}[Números con regla]
Si $2 \to 5$, $3 \to 10$, $4 \to 17$, ¿a qué va el 6?
\solucion $n^2 + 1$: $6 \to 37$.
\end{ejemplo}

\begin{trampa}
Asumir que el orden de las palabras es el mismo que en español. Muchos de estos idiomas inventados ponen el adjetivo
después (o antes); la correspondencia se halla por repeticiones, no por posición.
\end{trampa}

\begin{reto}
En un idioma, «lumo sapi» significa «perro grande», «tera sapi» significa «gato grande» y «tera nov» significa «gato
negro». ¿Cómo se dice «perro negro»?
\opciones{lumo nov}{nov lumo}{sapi nov}{lumo tera}
\end{reto}

\begin{retoresuelto}
«sapi» aparece en las dos primeras, como «grande»: sapi = grande. Entonces lumo = perro y tera = gato. De la tercera,
nov = negro. El orden es sustantivo y luego adjetivo: «lumo nov». \textbf{Opción A.}
\end{retoresuelto}

\taller
\piso{Entrena}{deducir}
\begin{ejercicios}
\item «bi la» = «agua fría», «la to» = «fría noche». ¿Qué es «to»? \sol{noche.}
\item Si $1 \to 3$, $2 \to 6$, $3 \to 11$, $4 \to 18$, ¿a qué va el 5? \sol{$n^2 + 2$: 27.}
\end{ejercicios}
\piso{Como en el examen}{opción múltiple}
\begin{ejercicios}
\item Si $4 \star 2 = 62$, $5 \star 3 = 82$ y $7 \star 1 = 86$, ¿cuánto es $6 \star 4$?
\opciones{102}{1\,024}{210}{224}
\sol{A. Se escribe la suma seguida de la diferencia: $4 + 2 = 6$ y $4 - 2 = 2$ dan 62; $5 + 3$ y $5 - 3$ dan 82; $7 + 1$ y $7 - 1$ dan 86. Para $6 \star 4$: 10 y 2, es decir, 102.}
\end{ejercicios}

\bitacora

\evaluacion
\begin{ejercicios}
\item Con César de $k = 3$, la palabra HOLA se escribe
\opciones{KROD}{JNKZ}{KRNA}{IPMB}
\sol{A. (7.1)}
\item Si 3-1-6-5 significa CAFE, ¿qué significa 4-1-4-15?
\opciones{DADO}{DEDO}{CACO}{DIDO}
\sol{A. (7.1)}
\item «mi ta» = «casa blanca», «ta po» = «blanca nube», «po su» = «nube gris». ¿Qué es «su»?
\opciones{gris}{nube}{casa}{blanca}
\sol{A. (7.2)}
\item Si $2 \to 8$, $3 \to 27$, $5 \to 125$, ¿a qué va el 4?
\opciones{16}{64}{12}{48}
\sol{B. Cubos. (7.2)}
\end{ejercicios}
